\documentclass[../../main.tex]{subfiles}

\begin{document}

\section{仿紧致性}\label{section:ch06-paracompactness}

仿紧致性概念是近年来出现的, 它是紧致性概念的一种最有用的推广, 特别是在拓扑学和微分几何中很有用. 这里, 我们只给出它的一个应用, 这就是在下一节将要证明的度量化定理.

我们熟悉的许多空间都是仿紧致的. 例如, 每一个紧致空间都是仿紧致的, 这一点可从定义直接推得. 每一个可度量化空间也是仿紧致的, 这就是我们将要证明的 A. H. Stone 推出的定理. 这样, 仿紧致空间就包含了我们已讨论过的两类重要空间, 同时, 它也包含其他的一些空间.

为了搞清楚仿紧致性是怎样由紧致性推广而来的, 我们回忆一下紧致性的定义: 一个空间 $X$ 称为紧致的, 如果 $X$ 的每一个开覆盖包含一个覆盖 $X$ 的有限子族. 等价地说就是:

\begin{quote}
一个空间 $X$ 是紧致的, 如果 $X$ 的任意一个开覆盖 $\mathcal{A}$ 有有限的开加细 $\mathcal{B}$ 覆盖 $X$.
\end{quote}

这个定义与通常的定义等价, 因为给定了这样的加细 $\mathcal{B}$ 之后, 我们就可以对 $\mathcal{B}$ 中每一个成员选取 $\mathcal{A}$ 的一个成员包含它, 于是就得到了 $\mathcal{A}$ 的一个覆盖 $X$ 的有限子族.

紧致性这种新的说法虽然比较繁琐, 但是它蕴涵着一种推广的契机.

\begin{definition}\label{definition:ch06-paracompact}
空间 $X$ 称为仿紧致的(paracompact), 如果 $X$ 的任意一个开覆盖 $\mathcal{A}$ 有一个覆盖 $X$ 的局部有限的开加细 $\mathcal{B}$.
\end{definition}

许多作者仿照 Bourbaki, 在仿紧致性的定义中也要求 $X$ 是一个 Hausdorff 空间, (Bourbaki 在紧致的定义中要求了 Hausdorff 条件)这里我们不遵循这一方式.

\begin{examples}
空间 $\mathbb{R}^\omega$ 是仿紧致的. 令 $X=\mathbb{R}^\omega$. 设 $\mathcal{A}$ 是 $X$ 的一个开覆盖. 令 $B_0=\varnothing$, 对于任意正整数 $m$, $B_m$ 表示球心是坐标原点、半径为 $m$ 的开球. 给定 $m$, 选取 $\mathcal{A}$ 的有限多个元素覆盖 $\overline{B_m}$ 并且都和开集 $X-\overline{B_{m-1}}$ 相交. 用 $\mathcal{C}_m$ 表示这有限多个开集族. 那么族 $\mathcal{C}=\bigcup\mathcal{C}_m$ 就是 $\mathcal{A}$ 的一个加细. 显然 $\mathcal{C}$ 是局部有限的. 因为开集 $B_m$ 只与其中有限多个元素相交(即 $\mathcal{C}_1\cup\cdots\cup\mathcal{C}_m$ 中的元素). 最后, $\mathcal{C}$ 覆盖 $X$. 因为对于任意给定 $x$, 设 $m$ 是满足 $x\in B_m$ 的最小整数, 根据定义可见, $x$ 属于 $\mathcal{C}_m$ 中的某一个元素.
\end{examples}

仿紧致空间的某些性质类似于紧致空间的某些性质. 例如, 仿紧致空间的子空间未必仿紧致, 但闭子空间一定是仿紧致的. 仿紧致的 Hausdorff 空间必定是正规的. 另一方面, 仿紧致空间也有与紧致空间不同的地方, 特别是两个仿紧致空间的积未必再是仿紧致的. 稍后, 我们将证明这一事实.

\begin{theorem}\label{theorem:ch06-paracompact-hausdorff-regular}
每一个仿紧致的 Hausdorff 空间 $X$ 是正规的.
\end{theorem}

\begin{proof}
这一证明类似于紧致的 Hausdorff 空间是正规空间的证明.

首先证明正则性. 设 $a$ 是空间 $X$ 的一个点, $B$ 是不包含点 $a$ 的一个闭子集. 根据 Hausdorff 条件, 对于 $B$ 中的任意点 $b$, 可以选取包含 $b$ 的一个开集 $U_b$, 使它的闭包仍然不包含点 $a$. 这些开集 $U_b$ 和开集 $X-B$ 一起构成了 $X$ 的一个覆盖, 取 $X$ 的一个局部有限开加细 $\mathcal{C}$ 覆盖 $X$. 将 $\mathcal{C}$ 中所有与 $B$ 相交的成员组成的子族记为 $\mathcal{D}$, 则 $\mathcal{D}$ 覆盖 $B$. 进而, 如果 $D\in\mathcal{D}$, 则 $\overline{D}$ 不包含点 $a$, 因为 $D$ 与 $B$ 相交, 所以它必定在某一个集合 $U_b$ 中, 而 $U_b$ 的闭包不包含 $a$ 点. 令
\begin{align*}
V=\bigcup_{D\in\mathcal{D}}D.
\end{align*}
则 $V$ 是 $X$ 中包含 $B$ 的开集. 又因 $\mathcal{D}$ 是局部有限的,
\begin{align*}
\overline{V}=\bigcup_{D\in\mathcal{D}}\overline{D},
\end{align*}
从而 $\overline{V}$ 也不包含点 $a$, 正则性得证.

为了证明正规性, 我们只须重复上面的论证, 用闭集 $A$ 代替点 $a$, 用正则性代替 Hausdorff 条件就行了.
\end{proof}

\begin{theorem}\label{theorem:ch06-paracompact-closed-subspace}
仿紧致空间的每一个闭子空间是仿紧致的.
\end{theorem}

\begin{proof}
设 $Y$ 是仿紧致空间 $X$ 的一个闭子空间, $\mathcal{A}$ 是由 $Y$ 中开集组成的 $Y$ 的一个覆盖. 对于每一个 $A\in\mathcal{A}$, 选取 $X$ 的一个开集 $A'$, 使得 $A'\cap Y=A$. 所有开集 $A'$ 连同开集 $X-Y$ 一起构成 $X$ 的一个覆盖. 同时, 设 $\mathcal{B}$ 是 $X$ 的局部有限开覆盖, 并且 $\mathcal{B}$ 加细这一覆盖, 于是, 族
\begin{align*}
\mathcal{C}=\{B\cap Y\mid B\in\mathcal{B}\}
\end{align*}
就是所要求的 $\mathcal{A}$ 的局部有限开加细.
\end{proof}

\begin{examples}
Hausdorff 空间 $X$ 的仿紧致子空间未必是 $X$ 的闭集. 事实上, 开区间 $(0,1)$ 是仿紧致的, 因为它同胚于 $\mathbb{R}$, 但它却不是 $\mathbb{R}$ 中的闭子集.
\end{examples}

\begin{examples}
仿紧致空间的子空间未必是仿紧致的. 空间 $\bar{S}_\Omega\times\bar{S}_\Omega$ 是紧致的, 因此是仿紧致的. 但是子空间 $S_\Omega\times\bar{S}_\Omega$ 却不是仿紧致的, 因为它是 Hausdorff 的, 却不是正规的.
\end{examples}

为了证明每一个可度量化空间都是仿紧致的这个重要定理, 我们还需要下面这一引理, 它由 E. Michael 给出, 并可用于其他目的.

\begin{lemma}\label{lemma:ch06-paracompact-equivalences}
设 $X$ 是一个正则空间. 则加在 $X$ 上的下列条件等价: $X$ 的每一个开覆盖有一个加细, 它是
\begin{enumerate}[(1)]
\item $X$ 的一个开覆盖, 并且是可数局部有限的.
\item $X$ 的一个覆盖, 并且是局部有限的.
\item $X$ 的一个闭覆盖, 并且是局部有限的.
\item $X$ 的一个开覆盖, 并且是局部有限的.
\end{enumerate}
\end{lemma}

\par
\begin{figure}[H]
\centering
\includegraphics[width=0.5\textwidth]{fig41-1.png}
\caption{}\label{figure:ch06-fig41-1}
\end{figure}

\begin{proof}
(4)$\Rightarrow$(1)是显然的, 对于我们打算证明的定理而言, 需要证明的是它的逆命题. 为此, 我们要通过(1)$\Rightarrow$(2)$\Rightarrow$(3)$\Rightarrow$(4)这样的步骤来做. 为方便起见, 我们已将相应的条件列在引理的条文中了.

(1)$\Rightarrow$(2). 设 $\mathcal{A}$ 是 $X$ 的一个开覆盖, $\mathcal{B}$ 是 $\mathcal{A}$ 的一个可数局部有限的开加细, 并且覆盖 $X$. 令
\begin{align*}
\mathcal{B}=\bigcup\mathcal{B}_n,
\end{align*}
其中每一个 $\mathcal{B}_n$ 都是局部有限的.

现在我们利用与以前用过的收缩法基本相同的方法, 从不同的 $\mathcal{B}_n$ 作出一些两两无交的集合. 给定 $i$, 令
\begin{align*}
V_i=\bigcup_{U\in\mathcal{B}_i}U.
\end{align*}
再对于每一个 $n\in\mathbb{Z}_+$ 和 $\mathcal{B}_n$ 的每一个成员 $U$, 定义
\begin{align*}
S_n(U)=U-\bigcup_{i<n}V_i.
\end{align*}
[注意 $S_n(U)$ 不一定是开的, 也不一定是闭的.] 令
\begin{align*}
\mathcal{C}_n=\{S_n(U)\mid U\in\mathcal{B}_n\}.
\end{align*}
则 $\mathcal{C}_n$ 是 $\mathcal{B}_n$ 的一个加细, 因为对于每一个 $U\in\mathcal{B}_n$, $S_n(U)\subset U$.

令 $\mathcal{C}=\bigcup\mathcal{C}_n$. 可以断言 $\mathcal{C}$ 就是所要求的 $\mathcal{A}$ 的局部有限加细, 并且覆盖 $X$.

设 $x$ 是 $X$ 的一个点, 我们要证明 $x$ 落入 $\mathcal{C}$ 的某一个元素中, 并且 $x$ 有一个邻域, 只与 $\mathcal{C}$ 中的有限多个元素相交. 考虑覆盖 $\mathcal{B}=\bigcup\mathcal{B}_n$, 令 $N$ 为满足以下条件的那个最小的整数: 在 $\mathcal{B}_N$ 中存在某元素包含着 $x$. 设 $U$ 是 $\mathcal{B}_N$ 中包含着 $x$ 的一个元素. 首先注意, 对于 $i<N$, $x$ 不属于 $\mathcal{B}_i$ 的任何元素, 所以 $x$ 在 $\mathcal{C}$ 的元素 $S_N(U)$ 中. 其次, 因为每一个族 $\mathcal{B}_n$ 是局部有限的, 所以对于每一个 $n=1,\cdots,N$ 选取 $x$ 的邻域 $W_n$, 使它仅和 $\mathcal{B}_n$ 中的有限多个成员相交. 如果 $W_n$ 与 $\mathcal{B}_n$ 的成员 $S_n(V)$ 相交, 那么它也必与 $\mathcal{B}_n$ 中的成员 $V$ 相交, 因为 $S_n(V)\subset V$. 因此 $W_n$ 仅和 $\mathcal{C}_n$ 的有限多个成员相交. 此外, 因为 $U$ 在 $\mathcal{B}_N$ 中, 所以对于 $n>N$, $U$ 不与 $\mathcal{C}_n$ 的任何成员相交. 于是 $x$ 的邻域
\begin{align*}
W_1\cap W_2\cap\cdots\cap W_N\cap U
\end{align*}
仅与 $\mathcal{C}$ 的有限多个成员相交.

(2)$\Rightarrow$(3). 设 $\mathcal{A}$ 是 $X$ 的一个开覆盖, $\mathcal{B}$ 是 $X$ 所有满足以下条件的开集 $U$ 构成的集族: $\overline{U}$ 包含在 $\mathcal{A}$ 的一个元素之中. 由正则性可见 $\mathcal{B}$ 覆盖 $X$. 再应用(2), 可以找到 $X$ 的一个覆盖 $\mathcal{C}$ 加细 $\mathcal{B}$, 并且 $\mathcal{C}$ 是局部有限的. 令
\begin{align*}
\mathcal{D}=\{\overline{C}\mid C\in\mathcal{C}\}.
\end{align*}
则 $\mathcal{D}$ 也覆盖 $X$. 根据\cref{lemma:ch06-locally-finite-properties}可见 $\mathcal{D}$ 是局部有限的, 并且加细 $\mathcal{A}$.

(3)$\Rightarrow$(4). 设 $\mathcal{A}$ 是 $X$ 的一个开覆盖. 根据(3), 选取 $\mathcal{B}$ 为 $X$ 的一个局部有限的覆盖并且加细 $\mathcal{A}$. (如果愿意的话, 也可以取 $\mathcal{B}$ 为闭加细, 但这不是必要的.) 我们希望将 $\mathcal{B}$ 的每一个元素 $B$ 稍加扩大, 使之成为一个开集. 此外, 还要把这种扩大做得充分小, 使得所得出的开集族仍然是局部有限的, 并且仍然是 $\mathcal{A}$ 的加细.

这一步包含着一个新的技巧. 前面多次采用的办法都是先按照某种方式将集合排序, 再用减掉前面所有集合的办法构造新集合. 那种办法是缩小集合, 而现在要扩大这些集合, 我们必须要有不同的做法. 我们将引进 $X$ 的一个局部有限闭覆盖 $\mathcal{C}$ 作为辅助, 然后再用它扩大 $\mathcal{B}$ 的成员.

对于 $X$ 的每一点 $x$, 存在 $x$ 的一个邻域仅与 $\mathcal{B}$ 的有限多个成员相交. 那么所有满足仅与 $\mathcal{B}$ 的有限多个成员相交这一条件的开集构成的族是 $X$ 的一个开覆盖. 再应用(3), 设 $\mathcal{C}$ 是这个覆盖的一个闭加细, 它覆盖 $X$, 并且是局部有限的. 于是 $\mathcal{C}$ 的每一个成员也仅与 $\mathcal{B}$ 的有限多个成员相交.

对于 $\mathcal{B}$ 的每一个成员 $B$, 令
\begin{align*}
\mathcal{C}(B)=\{C\mid C\in\mathcal{C}\text{ 和 }C\subset X-B\}.
\end{align*}
然后定义
\begin{align*}
E(B)=X-\bigcup_{C\in\mathcal{C}(B)}C.
\end{align*}
因为 $\mathcal{C}$ 是一个局部有限的闭集族, 根据\cref{lemma:ch06-locally-finite-properties}, $\mathcal{C}$ 的任意子族的成员之并仍然是闭集, 因此, 集合 $E(B)$ 是开集. 此外, 按定义, $B\subset E(B)$. (如\cref{figure:ch06-fig41-1}所示, 其中 $\mathcal{B}$ 的成员表示为闭圆域及线段, 而 $\mathcal{C}$ 的成员表示为闭正方形区域.)


现在我们可能把每一个 $B$ 扩得太大, 集族 $\{E(B)\}$ 可能不是 $\mathcal{A}$ 的加细了. 但这是容易弥补的. 对于每一个 $B\in\mathcal{B}$, 可以选取 $\mathcal{A}$ 中包含 $B$ 的成员 $F(B)$. 然后定义
\begin{align*}
\mathcal{D}=\{E(B)\cap F(B)\mid B\in\mathcal{B}\}.
\end{align*}
那么集族 $\mathcal{D}$ 便是 $\mathcal{A}$ 的一个加细. 因为 $B\subset(E(B)\cap F(B))$ 并且 $\mathcal{B}$ 覆盖 $X$, 所以 $\mathcal{D}$ 也覆盖 $X$.

最后, 我们证明 $\mathcal{D}$ 是局部有限的. 给定 $X$ 的一个点 $x$, 选取 $x$ 的邻域 $W$, 使得它只与 $\mathcal{C}$ 的有限多个成员, 比如说 $C_1,\cdots,C_k$ 相交. 下证 $W$ 只与 $\mathcal{D}$ 中有限多个元素相交. 因为 $\mathcal{C}$ 是 $X$ 的一个覆盖, 故 $W$ 被 $C_1,\cdots,C_k$ 所覆盖. 由此只需说明 $\mathcal{C}$ 中的每一个成员 $C$ 仅与 $\mathcal{D}$ 中有限多个成员相交. 若 $C$ 与集合 $E(B)\cap F(B)$ 相交, 则 $C$ 与 $E(B)$ 相交. 根据 $E(B)$ 的定义, $C$ 不可能含在 $X-B$ 中, 从而它必与 $B$ 相交. 因为 $C$ 仅与 $\mathcal{B}$ 中有限多个元素相交, 因此它最多与 $\mathcal{D}$ 中同样多个集合相交.
\end{proof}

\begin{theorem}\label{theorem:ch06-metrizable-paracompact}
每一个可度量化空间都是仿紧致的.
\end{theorem}

\begin{proof}
设 $X$ 是一个可度量化空间, 根据\cref{lemma:ch06-metrizable-countable-locally-finite-refinement}可见, $X$ 的每一个开覆盖 $\mathcal{A}$, 有 $X$ 的一个可数局部有限的开覆盖加细它, 前述引理蕴涵着开覆盖 $\mathcal{A}$ 有覆盖 $X$ 的一个局部有限的开加细.
\end{proof}

\begin{theorem}\label{theorem:ch06-regular-lindelof-paracompact}
每一个正则的 Lindelöf 空间是仿紧致的.
\end{theorem}

\begin{proof}
设 $X$ 是正则的 Lindelöf 空间. 给定 $X$ 的一个开覆盖 $\mathcal{A}$, 那么它有一个可数子族覆盖 $X$, 这个可数子族当然是可数局部有限的. 由前述引理可见, $\mathcal{A}$ 有局部有限的开加细覆盖 $X$.
\end{proof}

\begin{examples}
两个仿紧致空间的积未必是仿紧致的. 空间 $\mathbb{R}_l$ 是仿紧致的, 因为它是正则的 Lindelöf 空间. 但 $\mathbb{R}_l\times\mathbb{R}_l$ 却不是仿紧致的, 因为它是一个 Hausdorff 空间但不是正规的.
\end{examples}

\begin{examples}
空间 $\mathbb{R}^\omega$ 在积拓扑和一致拓扑下都是仿紧致的. 这是因为在这些拓扑下, $\mathbb{R}^\omega$ 都是可度量化的. 在箱拓扑下 $\mathbb{R}^\omega$ 是否是仿紧致的并不知道. (见第 32 节习题 5.)
\end{examples}

\begin{examples}
若 $J$ 是不可数的, 那么积空间 $\mathbb{R}^J$ 就不是仿紧致的. 因为 $\mathbb{R}^J$ 是 Hausdorff 的却不是正规的.
\end{examples}

仿紧致空间具有很多有用的性质, 其中一条就被用来处理空间 $X$ 上单位分拆的存在性. 对于这个概念, 其有限情形下的表达形式我们已在第 36 节中提到过, 现在我们来讨论更一般的情况. 先来回顾一个概念: 设 $\phi\colon X\to\mathbb{R}$, 那么 $\phi$ 的支撑就是所有满足 $\phi(x)\neq 0$ 的点 $x$ 构成的集合的闭包.

\begin{definition}\label{definition:ch06-partition-of-unity}
设 $\{U_\alpha\}_{\alpha\in J}$ 是 $X$ 的一个加标开覆盖. 一个加标的连续函数族
\begin{align*}
\phi_\alpha\colon X\longrightarrow[0,1]
\end{align*}
称为 $X$ 的由 $\{U_\alpha\}$ 所控制的一个单位分拆(partition of unity), 若:
\begin{enumerate}[(1)]
\item 对于每一个 $\alpha$, $\phi_\alpha$ 的支撑包含于 $U_\alpha$.
\item 加标族 $\{\phi_\alpha$ 的支撑$\}$ 是局部有限的.
\item 对于每一个 $x$, $\sum_\alpha\phi_\alpha(x)=1$.
\end{enumerate}
\end{definition}

条件(2)蕴涵着: 每一个 $x\in X$ 都有一个邻域, 除了有限多个 $\alpha$ 外, 其他的 $\phi_\alpha$ 在此邻域上取值都为 $0$. 因此我们可以明白条件(3)中``取和''的意思了. 我们解释这些是为了说明``和''是对那些不为 $0$ 的 $\phi_\alpha(x)$ 求取的.

我们现在来构造任意一个仿紧致的 Hausdorff 空间上的单位分拆. 正如在第 36 节中处理有限的情形那样, 我们先来证明一个``收缩引理''.

\begin{lemma}[收缩引理]\label{lemma:ch06-shrinking-lemma}
设 $X$ 是仿紧致的 Hausdorff 空间, $\{U_\alpha\}_{\alpha\in J}$ 是 $X$ 的加标开覆盖. 那么存在 $X$ 的一个局部有限的加标开覆盖 $\{V_\alpha\}_{\alpha\in J}$, 使得对于每一个 $\alpha$, 有 $\overline{V_\alpha}\subset U_\alpha$.
\end{lemma}

对于每一个 $\alpha$ 有 $\overline{V_\alpha}\subset U_\alpha$, 这一条件有时也被表述为: 集族 $\{\overline{V_\alpha}\}$ 是集族 $\{U_\alpha\}$ 的精加细(precise refinement).

\begin{proof}
令 $\mathcal{A}$ 为所有那些开集 $A$ 的族, 要求这些开集 $A$ 满足条件: $\overline{A}$ 包含于族 $\{U_\alpha\}$ 的某一个元素之中. 根据 $X$ 的正则性可见, $\mathcal{A}$ 是 $X$ 的一个覆盖. 因为 $X$ 是仿紧致的, 所以可以找到 $X$ 的一个局部有限的开覆盖 $\mathcal{B}$ 加细 $\mathcal{A}$. 我们用某一个指标集 $K$ 来一一地给 $\mathcal{B}$ 加标, 这样 $\mathcal{B}$ 中的成员就可以表示为 $B_\beta$, $\beta\in K$, 并且 $\{B_\beta\}_{\beta\in K}$ 是一个局部有限的加标族. 因为 $\mathcal{B}$ 加细 $\mathcal{A}$, 我们可以定义映射 $f\colon K\to J$, 使得对于每一个 $\beta\in K$, $f(\beta)\in J$, 使得
\begin{align*}
\overline{B_\beta}\subset U_{f(\beta)}.
\end{align*}
对于每一个 $\alpha\in J$, 我们定义 $V_\alpha$ 为集族
\begin{align*}
\mathcal{B}_\alpha=\{B_\beta\mid f(\beta)=\alpha\}
\end{align*}
中成员的并. (注意, 如果没有指标 $\beta$ 使得 $f(\beta)=\alpha$, $V_\alpha$ 便是空集.) 对于 $\mathcal{B}_\alpha$ 的每一个成员 $B_\beta$, 根据定义有 $\overline{B_\beta}\subset U_\alpha$. 因为集族 $\mathcal{B}_\alpha$ 是局部有限的, 故 $\overline{V_\alpha}$ 为集族 $\mathcal{B}_\alpha$ 中成员的闭包的并, 因此 $\overline{V_\alpha}\subset U_\alpha$.

最后, 我们来验证局部有限性. 给定 $x\in X$, 选取 $x$ 的一个邻域 $W$ 使得只有有限多个 $\beta$, 设为 $\beta=\beta_1,\cdots,\beta_k$, 使得 $B_\beta$ 与 $W$ 相交. 那么仅当 $\alpha$ 是指标 $f(\beta_1),\cdots,f(\beta_k)$ 中的某一个时, $W$ 与 $V_\alpha$ 相交.
\end{proof}

\begin{theorem}\label{theorem:ch06-partition-of-unity-existence}
设 $X$ 是仿紧致的 Hausdorff 空间, $\{U_\alpha\}_{\alpha\in J}$ 是 $X$ 的一个加标开覆盖. 则存在 $X$ 的被 $\{U_\alpha\}$ 控制的一个单位分拆.
\end{theorem}

\begin{proof}
我们先两次使用收缩引理, 以获得 $X$ 的一个局部有限的加标开覆盖 $\{W_\alpha\}$ 和 $\{V_\alpha\}$, 使得对于每一个 $\alpha$, 有 $\overline{W_\alpha}\subset V_\alpha$ 和 $\overline{V_\alpha}\subset U_\alpha$. 因为 $X$ 是正规的, 所以对于每一个 $\alpha$ 我们可以选取一个连续函数 $\psi_\alpha\colon X\to[0,1]$, 使得 $\psi_\alpha(\overline{W_\alpha})=\{1\}$ 并且 $\psi_\alpha(X-V_\alpha)=\{0\}$. 因为 $\psi_\alpha$ 的函数值非 $0$ 的点只在 $V_\alpha$ 中取到, 所以
\begin{align*}
(\psi_\alpha\text{ 的支撑})\subset\overline{V_\alpha}\subset U_\alpha.
\end{align*}
进而, 还可见 $\{\overline{V_\alpha}\}$ 是局部有限的. (因为一个开集只有当它与 $V_\alpha$ 相交的情况下, 它才能与 $\overline{V_\alpha}$ 相交.) 因此加标族 $\{\psi_\alpha$ 的支撑$\}$ 也是局部有限的. 注意 $\{W_\alpha\}$ 是 $X$ 的覆盖, 所以对于任意给定的 $x\in X$, 至少存在一个 $\psi_\alpha$ 在 $x$ 点上取正值.

无限和
\begin{align*}
\Psi(x)=\sum_\alpha\psi_\alpha(x)
\end{align*}
是有意义的. 因对任意一个 $x\in X$, 存在一个邻域 $W_x$, 使得在 $(\psi_\alpha$ 的支撑)中只有有限多个成员与之相交. 我们可以说这个无限和就是(有限多个)非 $0$ 项的和, 由此可见, 将 $\Psi$ 限制在 $W_x$ 上就等于有限多个连续函数的和. 因此它在 $X$ 上也是连续的, 并且总取正值. 现在定义
\begin{align*}
\phi_\alpha(x)=\psi_\alpha(x)/\Psi(x),
\end{align*}
这便得到了所求的单位分拆.
\end{proof}

数学中单位分拆常被用来``拼凑''函数, 即将局部定义的函数``拼凑''起来以获得一个全局有定义的函数. 第 36 节中对它们的使用已说明了这一过程, 这里是另一个例子.

\begin{theorem}\label{theorem:ch06-positive-function}
设 $X$ 是一个仿紧致的 Hausdorff 空间. $\mathcal{C}$ 是 $X$ 的一个子集族. 对于每一个 $C\in\mathcal{C}$, 设 $\varepsilon_C$ 是一个正数. 如果 $\mathcal{C}$ 是局部有限的, 那么存在一个连续函数 $f\colon X\to\mathbb{R}$, 使得对于所有 $x$ 有 $f(x)>0$, 同时对于 $x\in C$ 有 $f(x)\leqslant\varepsilon_C$.
\end{theorem}

\begin{proof}
选取 $X$ 的一个开覆盖, 它的每一个成员最多与 $\mathcal{C}$ 中有限多个元素相交. 将此开集族变成一个加标族 $\{U_\alpha\}_{\alpha\in J}$. 选取 $X$ 上的一个由 $\{U_\alpha\}$ 控制的单位分拆 $\{\phi_\alpha\}$. 给定 $\alpha$, 当 $C$ 取遍 $\mathcal{C}$ 中所有与 $\phi_\alpha$ 的支撑相交的元素时, 设 $\delta_\alpha$ 是所有 $\varepsilon_C$ 中极小的数. 若没有这样的元素 $C$, 令 $\delta_\alpha=1$. 然后定义
\begin{align*}
f(x)=\sum\delta_\alpha\phi_\alpha(x).
\end{align*}
因为所有的 $\delta_\alpha$ 都是正的, 因此 $f$ 也是正的. 下面我们要说明的是, 对于 $x\in C$, 有 $f(x)\leqslant\varepsilon_C$. 为此只需说明, 对于任意 $x\in C$ 和任意 $\alpha$, 有
\begin{align*}
\delta_\alpha\phi_\alpha(x)\leqslant\varepsilon_C\phi_\alpha(x);
\end{align*}
当 $\sum\phi_\alpha(x)=1$ 时, 所要求证的不等式成立. 若 $x\notin\phi_\alpha$ 的支撑, 则由于 $\phi_\alpha(x)=0$, 可见所求证的不等式是显然的. 若 $x\in\phi_\alpha$ 的支撑, 并且 $x\in C$, 那么 $C$ 与 $\phi_\alpha$ 的支撑相交, 因此 $\delta_\alpha\leqslant\varepsilon_C$. 所以上述不等式成立.
\end{proof}

\subsection*{习题}

\begin{enumerate}
\item 举例说明, 空间 $X$ 是仿紧致的, 但未必它的每一个开覆盖 $\mathcal{A}$ 有一个局部有限的子族覆盖 $\mathcal{A}$.
\item (a) 证明: 一个仿紧致的空间和一个紧致空间构成的积空间是仿紧致的. [提示: 使用管状引理.]
\begin{enumerate}[(b)]
\item 证明: $S_\Omega$ 不是仿紧致的.
\end{enumerate}
\item 每一个局部紧致的 Hausdorff 空间都是仿紧致的吗?
\item (a) 证明: 若 $X$ 具有离散拓扑, 那么 $X$ 是仿紧致的.
\begin{enumerate}[(b)]
\item 若 $f\colon X\to Y$ 是连续的, 并且 $X$ 是仿紧致的, 但 $Y$ 的子空间 $f(X)$ 未必是仿紧致的.
\end{enumerate}
\item 设 $X$ 是仿紧致的. 对于 $X$ 的任意加标开覆盖我们已经证明了``收缩引理'', 这里将要对 $X$ 的任意局部有限加标族证明``扩张引理''.

\textbf{引理} 设 $\{B_\alpha\}_{\alpha\in J}$ 是仿紧致的 Hausdorff 空间 $X$ 的子集构成的局部有限的加标族. 则存在 $X$ 的开子集构成的局部有限的加标族 $\{U_\alpha\}_{\alpha\in J}$, 使得对于每一个 $\alpha$, $B_\alpha\subset U_\alpha$.
\item (a) 设 $X$ 是正则空间. 若 $X$ 能表示成自身的可数个紧致子空间的并, 则 $X$ 是仿紧致的.
\begin{enumerate}[(b)]
\item 证明 $\mathbb{R}^\omega$ 作为 $\mathbb{R}^\omega$ 的子空间, 在箱拓扑下是仿紧致的.
\end{enumerate}
\item 设 $X$ 是正则空间.
\begin{enumerate}[(a)]
\item 若 $X$ 是有限个仿紧致的闭子空间的并, 那么 $X$ 是仿紧致的.
\item 若 $X$ 是可数个仿紧致的闭子空间的并, 并且这些闭子空间的内部覆盖 $X$, 那么 $X$ 是仿紧致的.
\end{enumerate}
\item 设 $p\colon X\to Y$ 是一个完备映射. (参见第 31 节的习题 7.)
\begin{enumerate}[(a)]
\item 证明: 若 $Y$ 是仿紧致的, 那么 $X$ 也是仿紧致的. [提示: 若 $\mathcal{A}$ 是 $X$ 的一个开覆盖, 选取由集合 $B$ 构成的 $Y$ 的一个局部有限开覆盖, 要求 $B$ 满足条件 $p^{-1}(B)$ 能被 $\mathcal{A}$ 中有限多个元素覆盖. 然后取 $p^{-1}(B)$ 与 $\mathcal{A}$ 中元素的交.]
\item 证明: 若 $X$ 是仿紧致的 Hausdorff 空间, 则 $Y$ 也是仿紧致的 Hausdorff 空间. [提示: 若 $\mathcal{B}$ 是 $X$ 的一个局部有限的闭覆盖, 则 $\{p(B)\mid B\in\mathcal{B}\}$ 是 $Y$ 的局部有限的闭覆盖.]
\end{enumerate}
\item 设 $G$ 是局部紧致的连通拓扑群. 证明 $G$ 是仿紧致的. [提示: 设 $U_1$ 是 $e$ 的有紧致闭包的一个邻域, 归纳地定义 $U_{n+1}=\overline{U_n}\cdot U_1$. 证明 $\overline{U_n}$ 的并在 $G$ 中是既开又闭的.] 没有连通的条件, 此结论依然成立, 但证明要困难些.
\item \textbf{定理} 设 $X$ 是局部紧致的和仿紧致的 Hausdorff 空间, 那么 $X$ 的每一个分支有一个可数基.

证明: 设 $X_0$ 是 $X$ 的一个分支, 那么 $X_0$ 是局部紧致的和仿紧致的. 设 $\mathcal{C}$ 是 $X_0$ 的一个局部有限的覆盖, 它的每一个成员是 $X_0$ 中的开集并且有紧致的闭包. 设 $U_1$ 是 $\mathcal{C}$ 中的一个非空成员, 一般地, $U_n$ 是 $\mathcal{C}$ 中所有与 $\overline{U_{n-1}}$ 相交的成员的并. 证明 $\overline{U_n}$ 是紧致的, 并且这些集合 $U_n$ 覆盖 $X_0$.
\end{enumerate}

\end{document}